\newpage
\subsection{Widerstandsbestimmung im Abgleichverfahren}\label{sec:1-2}

\subsubsection{Aufgabenstellung}

\begin{highlightblock}[Widerstandsbestimmung im Abgleichverfahren]
Beschreibung (Nicht direkt aus dem Skriptum kopieren!)
\end{highlightblock}

\subsubsection{Verwendetes Material}

\begin{table}[h]
	\centering
	\caption{Verwendetes Material}
	\begin{tabularx}{\textwidth}{XXX}
		\textbf{Material} & \textbf{Verwendung} & \textbf{Anmerkungen}\\ \midrule
		& & \\
		& & \\
		& & \\
	\end{tabularx}	
\end{table}

\subsubsection{Versuchsaufbau}

\textbf{Geräte Einstellungen}

\begin{minipage}[t]{.48\textwidth}
	Labornetzteil:
	\begin{itemize}[noitemsep]
		\item Spannung: 
		\item Strom: 
		\item Toleranz:
	\end{itemize}
\end{minipage}\hfill
\begin{minipage}[t]{.48\textwidth}
	Multimeter:
	\begin{itemize}[noitemsep]
		\item Einstellung: (Ohm, Volt, Ampere)
		\item Messbereich: 
		\item Garantiefehlergrenzen: 
	\end{itemize}
\end{minipage}

\textbf{Messschaltung}

\begin{figure}[h]
	\centering
	\begin{circuitikz}[scale=1.2]
		\draw (0,4) to[open, v=$U_H$, o-o] (0,0)
			-- ++(2, 0) coordinate (k1)
			to[vR=$R_2$, i<_=$I_2$, *-] ++(0,2) coordinate (k2)
			to[R=$R_1$, i_<=$I_1$, -*] ++(0,2) coordinate (k3)
			-- ++(-2,0);
		\draw (k1) -- ++(3,0)
			to[vR=$R_4$, i<_=$I_4$] ++(0,2) coordinate (k4)
			to[vR=$R_3$, i_<=$I_3$] ++(0,2) -- ++(-3,0);
		\draw (k2) to[ai, v^=$U_0$, i>^=$I_0$, *-*] (k4);
		\node[yshift=-20pt] at ($(k2)!0.5!(k4)$) {$R_0$};
	\end{circuitikz}
	\caption{Schaltung}
\end{figure}

Es sollen zwei Messungen je R2/R1 Verhältnis durchgeführt werden

\subsubsection{Berechnungen}

Abgleichbedingung: $R_1 R_4 = R_2 R_3$
\begin{itemize}
	\item R1 und R3 grob mit den werten aus \autoref{tab:brueckenabsch} Einstellen (Toleranzen)
	\item R4 fein einstellen und messen (Garantiefehler)
	\item R1 Berechnen
\end{itemize}

\[ R_1 = \frac{R_2R_3}{R_4} \]

\subsubsection{Ergebnisse}

Um Zeit zu sparen die Abgeglichenen Brücken in Schritt 1, 3 und 5 gleich mit dem Widerstand in \autoref{sec:1-3} Auslenken

\begin{table}[h]
	\centering
	\caption{Brückendaten}\label{tab:brueckendaten}
	\begin{tabularx}{\textwidth}{Y|YYYY|Y}
	Nummer & $R_2:R_1$ & $R_2$ / \si{\ohm} & $R_3$ / \si{\ohm} & $R_4$ / \si{\ohm} & $R_1$ / \si{\ohm} \\ \midrule
	\end{tabularx}
\end{table}

\subsubsection{Auswertung}

Messreihe statistisch auswerten. Am Ende vertrauensintervall angeben.

\subsubsection{Erkenntnisse}

